%%% The main file. It contains definitions of basic parameters and includes all other parts.

% Meta-data of your thesis (please edit)
\input metadata.tex

% Generate metadata in XMP format for use by the pdfx package
\input xmp.tex

%% Settings for single-side (simplex) printing
% Margins: left 40mm, right 25mm, top and bottom 25mm
% (but beware, LaTeX adds 1in implicitly)
%\documentclass[12pt,a4paper]{report}
%\setlength\textwidth{145mm}
%\setlength\textheight{247mm}
%\setlength\oddsidemargin{15mm}
%\setlength\evensidemargin{15mm}
%\setlength\topmargin{0mm}
%\setlength\headsep{0mm}
%\setlength\headheight{0mm}
% \openright makes the following text appear on a right-hand page
%\let\openright=\clearpage

%% Settings for two-sided (duplex) printing
% \documentclass[12pt,a4paper,twoside,openright]{report}
% \setlength\textwidth{145mm}
% \setlength\textheight{247mm}
% \setlength\oddsidemargin{14.2mm}
% \setlength\evensidemargin{0mm}
% \setlength\topmargin{0mm}
% \setlength\headsep{0mm}
% \setlength\headheight{0mm}
% \let\openright=\cleardoublepage

%% If the thesis has no printed version, symmetric margins look better
\documentclass[12pt,a4paper]{report}
\setlength\textwidth{145mm}
\setlength\textheight{247mm}
\setlength\oddsidemargin{10mm}
\setlength\evensidemargin{10mm}
\setlength\topmargin{0mm}
\setlength\headsep{0mm}
\setlength\headheight{0mm}
\let\openright=\clearpage

%% Generate PDF/A-2u
\usepackage[a-2u]{pdfx}

%% Prefer Latin Modern fonts
\usepackage{lmodern}
% If we are not using LuaTeX, we need to set up character encoding:
\usepackage{iftex}
\ifpdftex
\usepackage[utf8]{inputenc}
\usepackage[T1]{fontenc}
\usepackage{textcomp}
\fi

%% Further useful packages (included in most LaTeX distributions)
\usepackage{amsmath}        % extensions for typesetting of math
\usepackage{amsfonts}       % math fonts
\usepackage{amsthm}         % theorems, definitions, etc.
\usepackage{bm}             % boldface symbols (\bm)
\usepackage{booktabs}       % improved horizontal lines in tables
\usepackage{caption}        % custom captions of floating objects
\usepackage{dcolumn}        % improved alignment of table columns
\usepackage{floatrow}       % custom float environments
\usepackage{graphicx}       % embedding of pictures
\usepackage{indentfirst}    % indent the first paragraph of a chapter
\usepackage[nopatch=item]{microtype}   % micro-typographic refinement
\usepackage{paralist}       % improved enumerate and itemize
\usepackage[nottoc]{tocbibind} % makes sure that bibliography and the lists
			    % of figures/tables are included in the table
			    % of contents
\usepackage{xcolor}         % typesetting in color

% The hyperref package for clickable links in PDF and also for storing
% metadata to PDF (including the table of contents).
% Most settings are pre-set by the pdfx package.
\hypersetup{unicode}
\hypersetup{breaklinks=true}

% Packages for computer science theses
\usepackage{algpseudocode}  % part of algorithmicx package
\usepackage{algorithm}
\usepackage{fancyvrb}       % improved verbatim environment
\usepackage{listings}       % pretty-printer of source code

% You might want to use cleveref for references
% \usepackage{cleveref}

% Set up formatting of bibliography (references to literature)
% Details can be adjusted in macros.tex.
%
% BEWARE: Different fields of research and different university departments
% have their own customs regarding bibliography. Consult the bibliography
% format with your supervisor.
%
% The basic format according to the ISO 690 standard with numbered references
\usepackage[natbib,style=iso-numeric,sorting=none]{biblatex}
% ISO 690 with alphanumeric references (abbreviations of authors' names)
%\usepackage[natbib,style=iso-alphabetic]{biblatex}
% ISO 690 with references Author (year)
%\usepackage[natbib,style=iso-authoryear]{biblatex}
%
% Some fields of research prefer a simple format with numbered references
% (sorting=none tells that bibliography should be listed in citation order)
%\usepackage[natbib,style=numeric,sorting=none]{biblatex}
% Numbered references, but [1,2,3,4,5] is compressed to [1-5]
%\usepackage[natbib,style=numeric-comp,sorting=none]{biblatex}
% A simple format with alphanumeric references:
%\usepackage[natbib,style=alphabetic]{biblatex}

% Load the file with bibliography entries
\addbibresource{bibliography.bib}

% Our definitions of macros (see description inside)
\input macros.tex

%%% Title page and various mandatory informational pages
\begin{document}
%%% Title page of the thesis and other mandatory pages

%%% Inscriptions at the opening page of the hard cover

% We usually do not typeset the hard cover, but if you want to do it, change \iffalse to \iftrue
\iffalse

\pagestyle{empty}
\hypersetup{pageanchor=false}
\begin{center}

\large
Charles University

\medskip

Faculty of Mathematics and Physics

\vfill

{\huge\bf\ThesisTypeTitle}

\vfill

{\huge\bf\ThesisTitle\par}

\vfill
\vfill

\hbox to \hsize{\YearSubmitted\hfil \ThesisAuthor}

\end{center}

\newpage\openright
\setcounter{page}{1}

\fi

%%% Title page of the thesis

\pagestyle{empty}
\hypersetup{pageanchor=false}
\begin{center}

\centerline{\mbox{\includegraphics[width=166mm]{img/logo-en.pdf}}}

\vspace{-8mm}
\vfill

{\bf\Large\ThesisTypeTitle}

\vfill

{\LARGE\ThesisAuthor}

\vspace{15mm}

{\LARGE\bfseries\ThesisTitle\par}

\vfill

\Department

\vfill

{
\centerline{\vbox{\halign{\hbox to 0.45\hsize{\hfil #}&\hskip 0.5em\parbox[t]{0.45\hsize}{\raggedright #}\cr
Supervisor of the \ThesisTypeName{} thesis:&\Supervisor \cr
\ifx\ThesisType\TypeRig\else
\noalign{\vspace{2mm}}
Study programme:&\StudyProgramme \cr
\fi
}}}}

\vfill

Prague \YearSubmitted

\end{center}

\newpage

%%% A page with a solemn declaration to the thesis

\openright
\hypersetup{pageanchor=true}
\vglue 0pt plus 1fill

\noindent
I declare that I carried out this \ThesisTypeName{} thesis on my own, and only with the cited
sources, literature and other professional sources.
I understand that my work relates to the rights and obligations under the Act No.~121/2000 Sb.,
the Copyright Act, as amended, in particular the fact that the Charles
University has the right to conclude a license agreement on the use of this
work as a school work pursuant to Section 60 subsection 1 of the Copyright~Act.

\vspace{10mm}

\hbox{\hbox to 0.5\hsize{%
In \hbox to 6em{\dotfill} date \hbox to 6em{\dotfill}
\hss}\hbox to 0.5\hsize{\dotfill\quad}}
\smallskip
\hbox{\hbox to 0.5\hsize{}\hbox to 0.5\hsize{\hfil Author's signature\hfil}}

\vspace{20mm}
\newpage

%%% Dedication

\openright

\noindent
\Dedication

\newpage

%%% Mandatory information page of the thesis

\openright
{\InfoPageFont

\vtop to 0.5\vsize{
\setlength\parindent{0mm}
\setlength\parskip{5mm}

Title:
\ThesisTitle

Author:
\ThesisAuthor

\DeptType:
\Department

Supervisor:
\Supervisor, \SupervisorsDepartment

Abstract:
\Abstract

Keywords:
{\def\sep{\unskip, }\ThesisKeywords}

\vfil
}

% In Czech study programmes, it is mandatory to include Czech meta-data:

\ifx\StudyLanguage\LangCS

\vtop to 0.49\vsize{
\setlength\parindent{0mm}
\setlength\parskip{5mm}

Název práce:
\ThesisTitleCS

Autor:
\ThesisAuthor

\DeptTypeCS:
\DepartmentCS

Vedoucí bakalářské práce:
\Supervisor, \SupervisorsDepartment

Abstrakt:
\AbstractCS

Klíčová slova:
{\def\sep{\unskip, }\ThesisKeywordsCS}

\vfil
}

\fi

}

\newpage

%%% Further pages will be numbered, which LaTeX automatically enables at the next chapter start


%%% A page with automatically generated table of contents of the thesis

\tableofcontents

%%% Each chapter is kept in a separate file
\chapter{Introduction}

\noindent In computer science and software engineering, the idea of software verification is becoming more and more important. Verification checks whether the program or system operates correctly and fulfills the given properties. The goal is to detect bugs and errors during the development process.

One of the frameworks that is becoming popular is the Constrained Horn Clauses(CHC)\cite{10.1007/978-3-031-13185-1_2} framework. CHC is a fragment of First Order Logic modulo constraints that captures many program verification problems as constraint solving. The main advantage of CHC is that it separates modeling from solving by translating the program's behavior and properties into constrained language and then using a specialized CHC solver to solve various verification tasks across programming languages by deciding the satisfiability problem of a CHC system.

Golem\cite{blicha_golem_2023} is one such solver, which integrates the interpolating SMT solver OpenSMT\cite{10.1007/978-3-319-40970-2_35}. Golem currently implements six model-checking algorithms to solve the CHC satisfiability problem.

On top of solving the CHC satisfiability problem, each engine in Golem provides a validity witness for their answer. In software verification we can think of these witnesses as invariants for SAFE answers and a path to a counterexample for UNSAFE answer. By providing these witnesses, we can ensure that the engines answer is correct. Also to check, wheter the engine doesn't give false witnesses, Golem has built in an internal validator to verify the correctness of the witness. 

One such algorithm, that can be used to solve these problems is PDKind (Property-Directed K-induction)\cite{7886665}. PDKind is an IC3\cite{6148908}-based algorithm, that separates reachability checking and induction reasoning, allowing to replace the induction core with k-induction.

\section{Goals}
\noindent The goal of this work is to create a PDKind engine and integrate it into the Golem Horn Solver\cite{blicha_golem_2023}. 

We will start with defining the background around SMT solving. Then, we will introduce the Golem solver and analyze the PDKind algorithm to design a solution that integrates the PDKind algorithm as an engine into the Golem solver. Afterward, we will implement this design and add the generation of validity witnesses to the implementation. We will test the correctness and efficiency of the implementation using a set of benchmarks and compare the results with other engines. The correctness of the witnesses will be verified using a set of tests that utilize an internal function of the Golem solver, which can verify the correctness of the witnesses.


\chapter{Background}

\newtheorem{definition}{Definition}
\section{Transition System}
\noindent Transition systems provide a framework for modeling how systems change over time through state transitions. They are extensively used in formal methods, program analysis, and verification. Transition systems consist of a set of initial states, a set of states, and a transition relation that defines the evolution of states. Transition systems are used as an abstraction in computer science and software verification to solve several problems, including the Constrained Horn Clause's (CHC) satisfiability. By offering a formal method of representing the behavior of a computer system or an abstract computing device, they make it easier to analyze and verify properties. We will introduce the basic ideas of transition systems in this chapter, which will help us understand the following chapters.

\begin{definition}
    \textbf{Transition System \cite{7886665}:}
    A transition system is a 4-tuple ($S$, $\mathcal{F}$, $\mathcal{T}$, $\mathcal{I}$), where:
    \begin{itemize}
        \item \( S \) is a finite set of state variables \( \vec{x} \). Each \( x \in \vec{x} \) has a primed version \( x' \) representing its value in the next state. A \textbf{state} \( s \) is a valuation over \( \vec{x} \) that assigns a value \( s(x) \) to each \( x \in \vec{x} \).
        \item \( \mathcal{F} \) is a set of quantifier-free formulas, where each \( F(\vec{x}) \in \mathcal{F} \) is a \textbf{state formula} that holds in a state \( s \) if \( s \models F \).
        \item \( \mathcal{T} \) is a set of quantifier-free formulas, where each \( T(\vec{x}, \vec{x}') \in \mathcal{T} \) is a \textbf{transition formula} describing valid transitions between states. A transition from state \( s \) to \( s' \) is allowed if \( (s, s') \models T \) for some \( T \in \mathcal{T} \).
        \item \( \mathcal{I} \) is a set of formulas, where each \( I(\vec{x}) \in \mathcal{I} \) is a \textbf{state formula} describing initial states. A state \( s \) is an initial state if \( s \models I \).
    \end{itemize}
\end{definition}

\begin{definition}
    \textbf{Successor \cite{7886665}:}
    State $s$ has a successor $s'$, when $T(s(\vec{x}), s'(\vec{x}'))$ is true.
\end{definition}

\begin{definition}
    \textbf{k-Reachability \cite{7886665}:}
    State $s$ is k-reachable if there exists a sequence
    $I = s_0, s_1, s_2, \ldots, s_k = s$, where each $s_{i+1}$ is a successor of $s_i$.
\end{definition}

\section{Safety of Transition System}
\noindent An essential property in the analysis of transition systems is safety. We consider a transition system to be safe if it never reaches an unsafe state during its execution. This property is crucial in verifying the correctness of programs and systems modeled using transition systems.

\begin{definition}
    \textbf{Safety Property:}
    A property \( P \) is a \textbf{safety property} if it ensures that "nothing bad will ever happen" during the execution of the system. Formally, given a transition system \( (S, \mathcal{F}, \mathcal{T}, \mathcal{I}) \), a safety property \( P \) holds if for every reachable state \( s \in S \), the condition \( P(s) \) is satisfied.
\end{definition}

\begin{definition}
    \textbf{Invariant \cite{7886665}:}
    An \textbf{invariant} is a property that holds in all reachable states of a transition system. Formally, a property \( I \) is an invariant if, for a transition system \( (S, \mathcal{F}, \mathcal{T}, \mathcal{I}) \), every reachable state \( s \in S \) satisfies \( I(s) \).
\end{definition}

\noindent Safety properties can often be expressed as invariants, as they require that the system never reaches an unsafe state.

\vspace{\baselineskip}\noindent\textbf{Verifying Safety:} To verify whether a transition system satisfies a given safety property, we need to check whether all reachable states of the system satisfy the property. This can be done using different approaches:

\begin{itemize}
    \item \textbf{Explicit State Exploration:} Enumerate all reachable states and verify that each state satisfies $P$.
    \item \textbf{Bounded Model Checking (BMC):} Unroll the transition system up to a given bound $k$ and check whether an unsafe state is reachable within $k$ steps.
    \item \textbf{Inductive Invariant Checking:} Find an inductive invariant $Inv$ such that $\mathcal{I} \Rightarrow Inv$, $Inv \land \mathcal{T} \Rightarrow Inv'$, and $Inv \Rightarrow P$. If such an invariant exists, then $P$ holds for all reachable states.
\end{itemize}

\vspace{\baselineskip}\noindent\textbf{Example:} Consider a transition system modeling a counter with the following transitions:
\begin{align*}
    \mathcal{I}&: x = 0 \\
    \mathcal{T}&: x' = x + 1
\end{align*}

A safety property might state that the counter never exceeds a threshold, e.g., $x \leq 10$. The verification process would involve checking whether there exists a trace that leads to a state where $x > 10$. Which, in this situation, obviously exists.

Ensuring safety is critical for system reliability, and we can use various formal verification techniques to determine whether a system satisfies its safety properties.

\section{Satisfiability Modulo Theories}
\noindent In logic and computer science, the Boolean satisfiability problem, also called SAT, is one of the most well-known problems. In 1971, it was proven to be the first NP-complete problem \cite{10.1145/800157.805047}. This greatly helped in the field of complexity theory by providing a tool for the classification of other computational problems. SAT is widely used, but there are cases where it is not possible to represent a problem using only the two Boolean values, true and false.

This led to the generalization of SAT to more complex formulas, including real numbers, integers, or even lists and other data structures. In other words, an SMT instance arises as a generalization of a SAT instance, where Boolean variables are replaced with predicates of some theory. Examples of such theories are uninterpreted functions (UF), arrays, linear arithmetic (LA), or, more specifically, linear real arithmetic (LRA), which we will focus on in the rest of this work.

SMT solvers are widely used in various fields, such as formal verification, model checking, program synthesis, and hardware verification. They are instrumental in solving real-world problems that involve reasoning over complex theories, such as reasoning about hardware designs, program correctness, and even automated theorem proving.

The concept of SMT was introduced in the late 1990s as a way to extend SAT to handle more expressive logical formulas. SMT solvers combine techniques from both SAT solving and decision procedures for various theories, making them more powerful and versatile than traditional SAT solvers.

Each theory in SMT corresponds to a set of logical rules and operations for a specific type of data. For instance, in the case of Linear Real Arithmetic (LRA), the theory defines operations and constraints over real numbers. SMT solvers then check the satisfiability of a formula under a combination of such theories, making it more expressive than just Boolean logic.

For example, consider a system where we need to check if there is an assignment of integer values $x$ and $y$ that satisfies the following conditions:
\begin{equation*}
    2x + y \geq 10 \land x - y = 3
\end{equation*}

A traditional SAT solver would not be able to handle this problem, as it involves arithmetic over integers. However, using the linear arithmetic (LA) theory, an SMT solver can quickly solve this problem.

While SAT solvers are only used for assigning values for Boolean variables, SMT solvers work with more complex theories, allowing them to handle constraints over data types beyond just true or false values. This extension makes SMT a powerful tool for solving problems that SAT solvers cannot address.

\section{CHC Satisfiability}
\noindent Constrained Horn Clauses (CHCs) is a fragment of First Order Logic
modulo constraints that capture many program verification problems as
constraint solving.

\vspace{\baselineskip}\noindent
They extend Horn Clauses, which are logical formulas of the form:
\begin{equation*}
    P_{1} \land P_{2} \land \dots \land P_{n} \implies H
\end{equation*}
where $P_i$ are body literals and $H$ is the head of the clause.

\vspace{\baselineskip}
In CHCs, the literals in the body can include constraints in a first-order theory, such as linear arithmetics, boolean logic, etc.
Therefore, CHC is a formula of the form:
\begin{equation*}
    \phi \land P_{1} \land P_{2} \land \dots \land P_{n} \implies H
\end{equation*}
where $\phi$ is a constraint in the first order theory, $P_i$ and $H$ are uninterpreted predicates.

\vspace{\baselineskip}\noindent
A CHC is satisfiable if there exists an interpretation of the uninterpreted predciates $P_i$ and $H$ that is valid in the background theory.

The main advantage of CHC is that it separates the modeling
from solving by translating the program’s behavior and properties into
constrained language and then using a specialized CHC solver to solve various
verification tasks across programming languages by deciding the satisfiability
problem of a CHC system.

\vspace{\baselineskip}\noindent
For example, we consider a program with a loop:
\vspace{\baselineskip}
\begin{figure}[h]
    \begin{mdframed}
        \begin{lstlisting}
int i = 10;

while (x > 0) {
    x = x - 1;
}

assert (x >= 0);
        \end{lstlisting}
    \end{mdframed}
\end{figure}

\noindent We can encode the behavior and safety of this program in CHCs:

\noindent\vspace{\baselineskip}
\textbf{Initial clause:}
\begin{equation}
    x = 10 \implies P(x)
    \label{chc:init}
\end{equation}
\vspace{\baselineskip}
\textbf{Transition clause:}
\begin{equation}
    x > 0 \land P(x) \land x' = x - 1 \implies P(x')
    \label{chc:trans}
\end{equation}
\vspace{\baselineskip}
\textbf{Safety clause:}
\begin{equation}
    x < 0 \land P(x) \implies False
    \label{chc:safe}
\end{equation}

In this simple example, we encode the program using three CHC clauses. \eqref{chc:init} describes the initial state of the program and introduces an
uninterpreted predicate $P(x)$, which represents the loop invariant. \eqref{chc:trans} describes one step of the loop's transition: it checks if $x>0$
(the loop condition) and $P(x)$ (the invariant), then decrements $x$.
Satisfiability of the left-hand side of this clause implies that the invariant
$P(x')$ holds for the decremented value of $x$. Finally, \eqref{chc:safe}
checks the assertion condition. We verify the assertion by checking if its negation
leads to a contradiction. If the negation is satisfiable, it
implies a false state, indicating that the assertion does not hold.

In the previous example, we demonstrated how to encode a simple loop using
Constrained Horn Clauses (CHCs). This encoding process is not limited to simple
loops and can be extended to more complex and less deterministic loops. By
representing the loop’s initial state, transition, and termination conditions
as CHCs, we can analyze the program's behavior over a series of execution
steps. We can introduce a compact representation of a loop in CHCs as follows.

\vspace{\baselineskip}
\begin{equation}
    \begin{aligned}
        Init(V) &\implies Inv(V) \\
        Inv(V) \land Tr(V, V') &\implies Inv(V') \\
        Inv(V) \land Bad(V) &\implies False
    \end{aligned}
\end{equation}
This representation allows us to analyze the loop's behavior over a finite number of iterations. To obtain a finite trace of the loop, we unroll it $k$ times, where $k$ is a finite number.
\vspace{\baselineskip}
\begin{equation}
    Init(V) \land Tr(V, V^2) \dots \land Tr(V^{k-1}, V^k) \land Bad(V^k)
\end{equation}
This unrolled loop is known as a \textbf{Bounded Model Checking (BMC)} formula.
By increasing $k$, we can evaluate more BMC formulas, which helps detect more system bugs. However, we must identify inductive invariants to prove that no bugs exist. In our work, we will leverage BMC formulas to discover
these invariants.


\section{OpenSMT}
\noindent OpenSMT\cite{10.1007/978-3-642-12002-2_12} is an open-source SMT solver developed by the Formal Verification and Security Lab based at the University of Lugano. It includes a parser for SMT-LIB\cite{BarFT-SMTLIB} language, a state-of-the-art SAT-Solver, and a clean interface for new theory solvers, currently supporting various logics such as QF\_UF, QF\_LIA, QF\_LRA, and more. Currently, there is a second version OpenSMT2\cite{10.1007/978-3-319-40970-2_35}, which supports interpolation for propositional logic, QF\_UF, QF\_LRA, QF\_LIA, and QF\_IDL.

\section{Induction vs k-Induction}
\noindent The PDKind algorithm is a combination of IC3 and k-induction.
IC3\cite{6148908} is a commonly used method that uses induction to show a property is invariant by incrementally constructing an inductive strengthening of the property. PDKind \cite{7886665} breaks IC3 into modules, which allows replacing the induction method with k-induction.

We need to introduce two definitions to compare the relative strength of induction and k-induction. Let us have a transition system ($S$, $\mathcal{F}$, $\mathcal{T}$, $\mathcal{I}$).
\vspace{\baselineskip}
\begin{definition}
($\mathcal{F}$-Induction\cite{7886665}): Let $\mathcal{F}$ be a set of state formulas, where $P \in \mathcal{F}$. Then \( P \) is $\mathcal{F}$-inductive if:
\begin{align*}
    I(\vec{x})
    &\implies \mathcal{F}(\vec{x}), \hspace{5cm}&\text{(init)} \\
    \mathcal{F}(\vec{x}) \land T(\vec{x}, \vec{x}') &\implies P(\vec{x}'). \hspace{5cm}&\text{(cons)}
\end{align*}
If $\mathcal{F} = {P}$, we say that \( P \) is inductive,
\label{Def:ind}
\end{definition}

\vspace{\baselineskip}
\cite{7886665} states that if P is inductive, it is also invariant. However, invariants are generally not inductive. To prove that \( P \) is an invariant, we need to find a set of formulas $\mathcal{F}$ that includes \( P \) and is itself inductive. This set is also called the strengthening of \( P \). If such a strengthening exists, then \( P \) must be an invariant.

\begin{minipage}{\textwidth}
\begin{definition}
($\mathcal{F}^k$-Induction\cite{7886665}):Let $\mathcal{F}$ be a set of state formulas, where $P \in \mathcal{F}$. Then \( P \) is $\mathcal{F}^k$-inductive if:

\begin{align*}
    I(\vec{x_0}) \land \bigwedge_{i = 0}^{l - 1}T(\vec{x_i}, \vec{x}_{i+1})
    &\implies \mathcal{F}(\vec{x_l}), \quad &\text{for } 0 \leq l < k \quad \text{(k-init)} \\
    \bigwedge_{i = 0}^{k - 1}(\mathcal{F}(\vec{x_i}) \land T(\vec{x_i}, \vec{x}_{i+1}))
    &\implies P(\vec{x_k}). \quad &\text{(k-cons)}
\end{align*}
If $\mathcal{F} = {P}$, we say that \( P \) is \( k \)-inductive,
\label{Def:kind}
\end{definition}
\end{minipage}
If we substitute \( k \) with 1 in \( Definition \) \ref{Def:kind}, we can see that a property that is inductive is 1-inductive and also \( k \)-inductive for any \( k \). In the other direction, \cite{10.1007/978-3-319-08867-9_55} shows that if a property \( P \) is \( k \)-inductive and the logical theory allows quantifier elimination, then an inductive strengthening of \( P \) can be constructed by interpolation.

For theories that admit quantifier elimination, k-induction and induction have the same deductive power. However, \cite{Bjørner2015} shows that k-induction can yield more succinct strengthening. For other theories, such as \( Boolean logic \) or \( Linear arithmetic \), k-induction may be exponentially more powerful than induction due to weaker interpolation.

\subsection*{Practical Effectiveness}
\noindent \cite{7886665} shows that k-Induction is effective, especially when combined with algorithms like IC3.
\cite{7886665} demonstrates the effectiveness of k-induction on an example. The same example is depicted in Figure \ref{ex:ind vs kind}.

\begin{figure}[h]
    \begin{mdframed}
        \begin{lstlisting}
Invariant Property P: a[0] == 0

Initial States:
i = 0
j = 0
a[0] = 0

Transition:
j = Random() mod (i+1)
i = i + 1 mod N
a[i] = a[j]
        \end{lstlisting}
        \caption{Induction vs K-induction example}\label{ex:ind vs kind}
    \end{mdframed}
\end{figure}

\vspace{\baselineskip}The example shows a transition system, where zeroes are written to an array in a circular manner. The goal is to prove that the property $a[0] == 0$ is true throughout the execution of the program.

\newpage
\vspace{\baselineskip}\noindent \textbf{Induction:}

\vspace{\baselineskip}\noindent To prove \( P \) using induction, we must show the following:
\begin{itemize}
    \item\textbf{Initial state}: \( P \) holds in the initial state. This is true since $a[0] = 0$
    \item\textbf{Inductive step}: If \( P \) holds in state \( k \), it must hold in state \( k + 1\)
\end{itemize}

Here induction fails because it only considers a single step, which doesn't guarantee that $a[j]$ is always defined correctly. Thus, the program could copy a value from an undefined location to $a[0]$, which could break the property.

\vspace{\baselineskip}\noindent \textbf{K-Induction:}

\vspace{\baselineskip}\noindent While induction attempts to prove \( P \) in a single step, k-induction considers a sequence of steps. In this case, choosing \( k = N + 1 \) allows the index \( i \) to cycle back to 0, covering a complete cycle of transitions.

Over these \( N+1 \) steps, every position in the array copies a zero from a previous index, eventually filling the entire array with zeros. This guarantees that \( a[0] = 0 \) always holds.
Thus, \( P \) is \( (N+1) \)-inductive.

\vspace{\baselineskip}This example illustrates how \( k \)-induction succeeds where simple induction fails, particularly in systems with state transitions that unfold over multiple steps.

%\section{SMT-LIB}
%\noindent SMT-LIB\cite{BarFT-SMTLIB} is an international initiative that provides standard descriptions of background theories used in SMT solvers and develops common input and output languages for SMT solvers. It offers a unified framework for the description of SMT problems.  Additionally, SMT-LIB maintains a library of input problems, or benchmarks, written in the SMT-LIB language to support the development of different SMT solvers.


%\begin{figure}[h]
%\begin{mdframed}
%\begin{lstlisting}[language=Lisp]
%; Integer arithmetic
%(set-logic QF_LIA)
%(declare-const x Int)
%(declare-const y Int)

%; Assertion
%(assert (= (- x y) (+ x (- y) 1)))

%; Check satisfiability
%(check-sat)

%(exit)
%\end{lstlisting}
%\end{mdframed}
%\caption{SMT-LIB input example}
%\end{figure}

%In the Figure \ref{ex:ind vs kind} we can see an example of the SMT-LIB input taken from the offical examples of SMT-LIB\cite{BarFT-SMTLIB} .\begin{itemize}
%\item Lines starting with ';' represent comments and are ommited.
%\item \textbf{(set-logic QF\_LIA)}: This command sets the logic to Quantifier-Free Linear Integer Arithmetic (QF\_LIA), which involves linear integer arithmetic expressions without quantifiers.
%\item \textbf{(declare-const x Int) and (declare-const y Int)}: These commands declare two integer constants x and y.
%\item \textbf{(assert (= (- x y) (+ x (- y) 1)))}: This assertion states that $x - y$ should be equal to $x - y + 1$
%\item \textbf{(check-sat)}: This command asks the SMT solver to check the satisfiability of the assertion.
%\item \textbf{(exit)}: This command exits the solver.
%\end{itemize}

\chapter{Golem}

Golem \cite{blicha_golem_2023} is a flexible and efficient solver for CHC satisfiability problems over linear real arithmetic (LRA) and linear integer arithmetic (LIA), written in C++.

Golem integrates an interpolating SMT solver OpenSMT\cite{10.1007/978-3-319-40970-2_35}, and currently implements six different back-end engines for CHC solving, where each engine can use the OpenSMT not only for SMT queries but also for interpolant computation.

\section{Solver overview}

\begin{figure}[ht]
    \centering
\begin{tikzpicture}[node distance=1.5cm]

% Nodes
\node (smt2) [io] {.smt2};
\node (interpreter) [process, right of=smt2, xshift=1.2cm] {Interpreter};
\node (preprocessor) [process, right of=interpreter, xshift=1.2cm] {Preprocessor};
\node (engines) [process, right of=preprocessor, xshift=1.2cm] {Engines};
\node (opensmt) [decision, right of=engines, xshift=1.5cm] {OpenSMT};
\node (result) [result, below of=interpreter] {Result};
\node (sat) [below of=preprocessor, xshift=0.5cm, yshift=0.5cm, font=\footnotesize] {\textcolor{OliveGreen}{SAT + Model}};
\node (unsat) [below of=preprocessor, xshift=0.7cm, yshift=-0.5cm, font=\footnotesize] {\textcolor{red}{UNSAT + Proof}};
% Arrows
\draw [arrow] (smt2) -- (interpreter);
\draw [arrow] (interpreter) -- (preprocessor);
\draw [arrow] (preprocessor) -- (engines);
\draw [arrow] (engines) -- (opensmt);
\draw [arrow] (opensmt) -- (engines);
\draw [arrow, color=OliveGreen] (result) -- (sat);
\draw [arrow, color=red] (result) -- (unsat);

\draw [thick,->,>=stealth, color=blue] (engines.north) to [out=90,in=90] (preprocessor.north);
\draw [thick,->,>=stealth, color=blue] (preprocessor.north) to [out=90,in=90] (interpreter.north);
\draw [thick,->,>=stealth, color=blue] (interpreter) -- (result);
\end{tikzpicture}
\caption{Architecture of Golem}
    \label{fig:golem_diagram}
\end{figure}


\noindent In this section, we will describe the Golem solving process depicted in Figure \ref{fig:golem_diagram}.
\section*{Reading and interpreting CHCs}
\noindent Golem reads the input from a file in the .smt2 format, which is an extension of the SMT-LIB language\cite{BarFT-SMTLIB}. The interpreter builds an internal representation of the CHC system by first normalizing the CHCs to ensure that each predicate has only variables as arguments and then converting the CHCs to the graph representation. The graph representation is then passed to the preprocessor.

\section*{Preprocessing}
\noindent The Preprocessor applies transformations to simplify the graph representation:
        \begin{itemize}
            \item \textbf{Predicate Elimination:} Removes predicates, that are not present in both the body and the head of the same clause.
            \item \textbf{Clause Merging:} Merges clauses with the same uninterpreted predicate by disjoining their constraints.
            \item \textbf{Redundant Clause Deletion:} Removes clauses, that cannot participate in the proof of unsatisfiability. 
        \end{itemize}
    \section* {Engines:}
    \noindent The graph is then solved with one of the engines (this option is specified by the user):
        \begin{itemize}
            \item Bounded Model Checking (BMC) \cite{10.1007/3-540-49059-0_14}
            \item k-Induction (KIND) \cite{10.1007/3-540-40922-X_8}
            \item Interpolation-based Model Checking (IMC) \cite{10.1007/978-3-540-45069-6_1}
            \item Lazy Abstractions with Interpolants (LAWI) \cite{10.1007/11817963_14}
            \item Spacer \cite{10.1007/978-3-319-08867-9_2}
            \item Transition Power Abstraction (TPA) \cite{blicha_golem_2023}
        \end{itemize}

The user can also select the option to produce a validity witness. When the engine solves the problem, it generates a model for the SAT result or a proof for the UNSAT result. These models and proofs are translated back by the preprocessor to match the original system.

\begin{figure}[H]
\begin{align*}
    x \leq 1 &\implies I(x) \quad &\text{\color{red}{\(I(1)\)}} \\
    x' = x + 1 &\implies T(x, x') \quad &\text{\color{red}{\(T(1, 2)\)}} \\
    I(x) \land T(x, x') &\implies S(x') \quad &\text{\color{red}{\(S(2)\)}} \\
    S(x) \land x \geq 2 &\implies \text{\color{red}{false}} \quad &\text{\color{red}{false}}
\end{align*}
    \caption{UNSAT + Proof example}\label{ex:UNSAT}
\end{figure}

In Figure \ref{ex:UNSAT}, we have a CHC system and a proof of its unsatisfiability. There are four derivation steps. The first step sets the variable $x$ to $1$ and gets I(1). The second step sets $x := 1$ and $x' := 2$ and gets T(1, 2). Step three applies resolution to the instance of the third clause for $x := 1$ and $x' := 2$ and the previously derived facts L(1) and T(1,2) and gets S(2). The last step again applies resolution to the instance of the fourth clause with $x := 2$ and the derived fact S(2) and gets false. 


\section{Engine integration}
\noindent The architecture described above allows us to integrate an engine into the Golem solver without modifying it. There are several ways we can do it. 

Writing a library in a different programming language could be interesting because we could leverage the pros of other languages. The downside is that the interaction of C++ and another language could bring various overheads. Our goal is to create an efficient engine, and therefore, having these overheads would not be optimal.

Another option would be to write a C++ library and include it in the project. This would eliminate the downsides of the first option, but there would still be some issues. For example, we would have to integrate an SMT solver for satisfiability checking and interpolation. 

The most efficient option is to integrate the engine directly into the solver the same way the other engines are integrated. This would allow us to directly utilize OpenSMT, which is integrated into Golem, for satisfiability checking and interpolation. We could also use many other features that Golem has, making our development easier.



\chapter{PDKind}

% Describe the individual procedures of the algorithm and in each part analyze how the implementation went. If there were more ways to do it, compare them.
% The last section should describe how we also added validity witnesses because Golem and other engines in it produce them).
\noindent PDKind is an IC3\cite{6148908}-based algorithm, that separates reachability checking and induction reasoning, allowing the induction core to be replaced with k-induction.

In this chapter, we describe and analyze the PDKind algorithm. In sections 4.1-4.3, we will be using methodology, algorithms, and definitions described in \cite{7886665}.
After reading this chapter, the reader should be able to understand the parts of the PDKind algorithm and orientate in our implementation.
%We will also describe the modifications of the algorithm and the decisions we made when implementing some abstract data structures from the paper.




\section{Reachability checking procedure}\label{ReachabilityChecking}

\noindent To understand the following, we need to introduce a few definitions.

\vspace{\baselineskip}\noindent Let us have a formula in the form

\begin{equation}
    A(\vec{x}) \wedge T[B]^k(\vec{x},\vec{y},\vec{w}) \wedge C(\vec{y})
\end{equation}

\vspace{\baselineskip}\begin{definition}
                          For \( k > 1 \), \( T[F]^k(\vec{x}, \vec{x}') \) is defined as
                          \begin{equation*}
                              T(\vec{x}, \vec{w}_1) \land \bigwedge_{i=1}^{k-2} \left( F(\vec{w}_i) \land T(\vec{w}_i, \vec{w}_{i+1}) \right) \land F(\vec{w}_{k-1}) \land T(\vec{w}_{k-1}, \vec{x}')
                          \end{equation*}

                          \noindent where $\vec{w}$ are state variables in the intermediate states.
\end{definition}


\vspace{\baselineskip}
\begin{definition}{\textbf{(Interpolant)}}
    \cite{7886665} If formula (4.1) is unsatisfiable, then $I(\vec{y})$ is an interpolant if
    \label{interpolation}
    \begin{enumerate}
        \item \( A(\vec{x}) \land T[B]^k(\vec{x}, \vec{w}, \vec{y}) \Rightarrow I(\vec{y}) \), and
        \item \( I(\vec{y}) \) and \( C(\vec{y}) \) are inconsistent.
    \end{enumerate}
\end{definition}

\vspace{\baselineskip}
\begin{definition}{\textbf{(Generalization)}}
    \cite{7886665} If formula (4.1) is satisfiable, then $G(\vec{x})$ is a generalization if
    \label{generalization}
    \begin{enumerate}
        \item $G(\vec{x}) \Rightarrow \exists \vec{y}, \vec{w}$  $T[B]^k(\vec{x}, \vec{y}, \vec{w}) \wedge C(\vec{y})$, and
        \item $G(\vec{x})$ and $A(\vec{x})$ are consistent.
    \end{enumerate}
\end{definition}

\renewcommand{\figurename}{Algorithm}
\begin{figure}[h]
    \begin{mdframed}
        \begin{algorithmic}[1]
            \State \textbf{Input:} Target state $F$, maximum steps $k$
            \State \textbf{Data:} Reachability frames $R$, initial states $I$, transition states $T$
            \State \textbf{Output:} True if $F$ is reachable in $k$ steps, False otherwise

            \If{$k = 0$}
                \State \Return{CheckSAT($I \wedge T^0 \wedge F$)}
            \EndIf

            \While{true}
                \If{CheckSAT($R_{k-1} \wedge T \wedge F$)}
                    \State $G \gets$ Generalize($R_{k-1}, T, F$)
                    \If{Reachable($G, k-1$)}
                        \State \Return{true}
                    \Else
                        \State $E \gets$ Explain($G, k-1$)
                        \State $R_{k-1} \gets R_{k-1} \cup E$
                    \EndIf
                \Else
                    \State \Return{false}
                \EndIf
            \EndWhile

        \end{algorithmic}
    \end{mdframed}
    \caption{Reachable method}\label{alg:1}
\end{figure}

\vspace{\baselineskip}
\noindent Method shown in Algorithm \ref{alg:1} tries to reach the initial states backwards by using a depth-first search strategy.

To check if \( F \) is reachable from the initial states in \( k \) steps, we
first check whether \( F \) is reachable in one transition from the previous
frame \( R_{k-1} \). If there is no such transition, then \( F \) is not
reachable in \( k \) steps. Otherwise, we get a state that satisfies \( R_{k-1}
\) and from which \( F \) is reachable in a single step. 

We then call a
generalization procedure, which produces a formula \( G \) (see Definition
\ref{generalization}). Here, generalization helps avoid redundant checks by creating an abstraction of the counterexamples. Instead of storing the given state, it generates a formula covering multiple possible counterexamples, which reduces the number of times the reachability procedure gets called.

Then, using a DFS strategy, we recursively
check whether \( G \) is reachable from the initial states.
If \( G \) is reachable, then \( F \) is also reachable, and the procedure terminates. Otherwise,
we can learn an explanation through interpolation (see Definition
\ref{interpolation}) and eliminate \( G \) by adding the
explanation into the frame \( R_{k-1} \).


\section{Push procedure} \label{Push}

\begin{definition}{\textbf{(Induction Frame)}}
    \cite{7886665}\label{def:IFrame} A set of tuples $F \subset \mathbb{F} \times \mathbb{F}$, where $\mathbb{F}$ is a set of all state formulas in theory $T$, is an induction frame at index $n$ if $(P, \neg P) \in F$ and $\forall (lemma, counterExample) \in F$:
    \begin{itemize}
        \item $lemma$ is valid up to $n$ steps and refutes $counterExample$
        \item $counterExample$ states can be extended to a counterexample to $P$.
    \end{itemize}
\end{definition}

\vspace{\baselineskip}
\noindent Induction frame tracks the strengthening of invariants over number of iterations. Every frame consists of a formula that has been refined to eliminate an invalid state.

The procedure shown in Algorithm \ref{alg:3} is the core of the PDKind algorithm.
It takes the invariant and tries to construct a stronger inductive proof by analyzing its counterexamples and strengthening constraints, while switching between \\
generalization and reachability checking.
We will break it down into smaller pieces to understand how it works.


The first step (line 11) is to check whether \( lemma \) is
\( k \)-inductive (see Definition \ref{Def:kind}). This is done by checking if the formula $(F_{ABS} \wedge T_k
\wedge \neg lemma)$ is satisfiable. We can notice, that \( CheckSAT \) call returns model $m_1$ which
is an assignment of values to variables that satisfies the formula.
If the formula $(F_{ABS} \wedge T_k \wedge \neg lemma)$ isn't satisfiable, we
can push a new obligation to our new induction frame \( G \) and continue.
Otherwise, we get the model $m_1$ and save it for later.

In the second part (line 16), we check if the \( counterExample \) is
reachable. If it is, we get model $m_2$ and generalize it to
$g_2$. We know, that we can reach \( \neg P \) from $g_2$,
so we need to check if $g_2$ is reachable from the initial states. If it is
reachable, the property is invalid, and we mark \( isInvalid \gets true \).


On line 19, we can see that \( Reachable() \) accepts more arguments and
returns more values than shown in Algorithm
\ref{alg:1}. This new
$Reachable(i,j,F)$ method checks if $F$ is reachable in \( k \) steps
where $i\leq k \leq j$.
If this call wasn't successful, we get an interpolant $i_1$ and assign
it to $g_3$ which eliminates $g_2$. We found a new induction
obligation $(g_3, g_2)$, which is a strengthening of \( F \). Now,
we can try again with a potential counterexample eliminated.

Last step is to analyze the induction failure. From the first check we have a
model $m_1$, which is a counterexample to the k-inductiveness of
\( lemma \). We again get $g_1$ as a generalization of $m_1$
and check if $g_1$ is reachable from initial states. If it is
reachable, we replace \( lemma \) with weaker \( \neg counterExample \)
and push this new obligation to \( F \) and \( G \). On the other hand,
if $g_1$ is not reachable, we strengthen \( lemma \) with
$g_3$ and push this new obligation to \( F \).

\begin{figure}[H]
    \begin{mdframed}
        \begin{algorithmic}[1]
            \State \textbf{Input:} Induction frame $F$, $n$, $k$
            \State \textbf{Output:} Old induction frame $F$, new induction frame $G$, $n_p$, $isInvalid$

            \State push elements of $F$ to queue $Q$
            \State $G \gets$ \{\}
            \State $n_p \gets n + k$
            \State $invalid \gets false$
            \While {$\neg invalid$ and $Q$ is not empty}
                \State $(lemma, counterExample)  \gets$ $Q$.pop()
                \State $F_{ABS} \gets \bigwedge a_i$ ,where ($a_i$, $b_i$) \in $F$ \forall $i$ \in \{1,...,$F$.length()\}
                \State $T_k \gets T[F_{ABS}]^k$ by definition
                \State $(s_1, m_1) \gets$ CheckSAT($F_{ABS}, T_k, \neg lemma$)
                \If{$\neg s_1$}
                    \State $G \gets G \cup (lemma, counterExample)$
                    \State Continue
                \EndIf
                \State $(s_2, m_2) \gets$ CheckSAT($F_{ABS} \wedge T_k \wedge counterExample$)
                \If{$s_2$}
                    \State $ g_2 \gets $ Generalize($m_2, T_k, counterExample$)
                    \State $(r_1, i_1, n_1) \gets $ Reachable($n-k+1$, $n$, $g_2$)
                    \If{$r_1$}
                        \State $ isInvalid \gets true $
                        \State Continue
                    \Else
                        \State $g_3 \gets i_1$
                        \State $F \gets F \cup (g_3, g_2)$
                        \State $Q$.push($(g_3, g_2)$)
                        \State $Q$.push($lemma, counterExample$)
                        \State Continue
                    \EndIf
                \EndIf
                \State $g_1 \gets$ Generalize($m_1$, $T_k$, $\neg lemma$)
                \State $(r_2, i_2, n_2) \gets $ Reachable($n-k+1$, $n$, $g_1$)
                \If{$r_2$}
                    \State $(r_3, i_3, n_3) \gets $ Reachable($n+1$, $n_2 + k$, $g_1$)
                    \State $n_p \gets $ Min($n_p$, $n_3$)
                    \State $F \gets$ $F$ $\cup $ ($\neg counterExample$, $counterExample$)
                    \State $G \gets$ $G$ $\cup $ ($\neg counterExample$, $counterExample$)
                \Else
                    \State $g_3 \gets$ $i_2 \wedge lemma$
                    \State $F \gets$ $F$ $\cup $ ($g_3$, $counterExample$)
                    \State $F \gets$ $F$ $\setminus $ ($lemma$, $counterExample$)
                    \State $Q$.push(($g_3$, $counterExample$))
                \EndIf
                \State \Return ($F$, $G$, $n_p$, $isInvalid$)
            \EndWhile

        \end{algorithmic}
    \end{mdframed}
    \caption{Push procedure}\label{alg:3}
\end{figure}

\newpage
\section{PD-Kind procedure}

\noindent The main PDKind procedure shown in Algorithm \ref{alg:4}
checks if property \( P \) is invariant by
iteratively calling the \( Push \) procedure to find a k-inductive
strengthening of \( P \) for some $1 \leq k \leq n + 1$.
PDKind iteratively strengthens the invariant until it becomes inductive.
The strengthening
\( G \) is k-inductive, and if $F == G$, then no more strengthening is required, and \( P \) is proven to be invariant
,so we return \( SAFE \). If the \( Push \) procedure marks
\( isInvalid \) as \( True \), then a counterexample has been found, proving that the property is not invariant and we
return \( UNSAFE \). Otherwise, we update \( n \) and repeat the loop.


\begin{figure}[H]
\begin{mdframed}
\begin{algorithmic}[1]
\State \textbf{Input:} Initial states $I$, transition formula $T$, property $P$
\State \textbf{Output:} Return UNSAFE if P is invalid or SAFE when there is no inductive strengthening left

\State $n \gets 0$
    \State $F \gets (P, \neg P)$
    \While{true}
        \State $k \gets n + 1$
        \State $(F, G, n_p, isInvalid) \gets$ Push($F$, $n$, $k$)
        \If{$isInvalid$}
        \State \Return UNSAFE
        \EndIf
        \If{$F = G$}
        \State \Return SAFE
        \EndIf
        \State $n \gets n_p$
        \State $F \gets G$
    \EndWhile

\end{algorithmic}
\end{mdframed}
\caption{Main PD-Kind procedure}\label{alg:4}
\end{figure}


\chapter{Implementation}
% Zde bych popsal pouzity programovaci jazyk, doporucenou platformu.

% Nevim jestli tato cast nebude moc kratka, ale popsal bych jednotlive, jak reprezentuji ruzne datove struktury popsane v analyze.
\section{Data structures}

\subsection{Reachability frame}

\subsection{Induction frame}

% Zde bych popsal rozdeleni kodu na tridu, ktera se stara o Reachability check a na engine PDKIND
\section{Classes}

\subsection{Reachability class}

\subsection{PD-Kind engine}

\chapter{Experiments}

\chapter{Conclusion}
During this work, we introduced ourselves with the concept of the CHC framework. Then, we defined the concepts of transition systems and satisfiability modulo theories. After that, we described the structure of Golem and analyzed the best way to integrate a new engine into it. We then analyzed the entire PDKind algorithm and modified it to serve our needs.

With this knowledge, we implemented the PDKind engine itself. The goal was to create an engine that would 1) be well integrated into the Golem framework, 2) return correct answers, and 3) match the performance of other engines. To validate these goals, we experimentally compared our engine with other existing Golem engines. We managed to achieve results comparable to other engines, and in some cases, our engine was faster, although it did not reach the performance of the best-performing engine.

We believe that this engine presents a well-performing alternative to other engines, with its usefulness expected to grow with further improvements.


%%% Bibliography
%%% Bibliography (literature used as a source)
%%%
%%% We employ biblatex to construct the bibliography. It processes
%%% citations in the text (e.g., the \cite{...} macro) and looks up
%%% relevant entries in the bibliography.bib file.
%%%
%%% See also biblatex settings in thesis.tex.

%%% Generate the bibliography. Beware that if you cited no works,
%%% the empty list will be omitted completely.

% We let bibliography items stick out of the right margin a little
\def\bibfont{\hfuzz=2pt}

\printbibliography[heading=bibintoc]

%%% If case you prefer to write the bibliography manually (without biblatex),
%%% you can use the following. Please follow the ISO 690 standard and
%%% citation conventions of your field of research.

% \begin{thebibliography}{99}
%
% \bibitem{lamport94}
%   {\sc Lamport,} Leslie.
%   \emph{\LaTeX: A Document Preparation System}.
%   2nd edition.
%   Massachusetts: Addison Wesley, 1994.
%   ISBN 0-201-52983-1.
%
% \end{thebibliography}


%%% Figures used in the thesis (consider if this is needed)
\listoffigures

%%% Tables used in the thesis (consider if this is needed)
%%% In mathematical theses, it could be better to move the list of tables to the beginning of the thesis.
\listoftables

%%% Abbreviations used in the thesis, if any, including their explanation
%%% In mathematical theses, it could be better to move the list of abbreviations to the beginning of the thesis.
\chapwithtoc{List of Abbreviations}

%%% Doctoral theses must contain a list of author's publications
\ifx\ThesisType\TypePhD
\chapwithtoc{List of Publications}
\fi

%%% Attachments to the thesis, if any. Each attachment must be referred to
%%% at least once from the text of the thesis. Attachments are numbered.
%%%
%%% The printed version should preferably contain attachments, which can be
%%% read (additional tables and charts, supplementary text, examples of
%%% program output, etc.). The electronic version is more suited for attachments
%%% which will likely be used in an electronic form rather than read (program
%%% source code, data files, interactive charts, etc.). Electronic attachments
%%% should be uploaded to SIS. Allowed file formats are specified in provision
%%% of the rector no. 72/2017. Exceptions can be approved by faculty's coordinator.
\appendix
\chapter{Attachments}

\section{First Attachment}

\end{document}
